\documentclass[11pt]{article}
\usepackage[margin=0.82in]{geometry}
\usepackage{amsmath,amssymb,mathtools}
\usepackage{booktabs,tabularx}
\usepackage[dvipsnames]{xcolor}
\usepackage[colorlinks=true,linkcolor=NavyBlue,urlcolor=BrickRed,hypertexnames=false]{hyperref}
\usepackage{microtype}
\usepackage{fancyhdr}
\usepackage{enumitem}

\definecolor{ink}{HTML}{211B32}
\definecolor{accent}{HTML}{6D3B72}
\definecolor{pale}{HTML}{F6F0F7}
\definecolor{passgreen}{HTML}{16794A}
\color{ink}
\pagestyle{fancy}
\fancyhf{}
\lhead{$H_3$ reflection group}
\rhead{Proof and verification}
\cfoot{\thepage}
\setlength{\headheight}{14pt}
\newcommand{\Sfun}{\mathcal S}
\newcommand{\Sym}{\operatorname{Sym}}
\newcommand{\PASS}{\textcolor{passgreen}{\textbf{PASS}}}
\newcommand{\summarybox}[1]{\begin{center}\fcolorbox{accent}{pale}{\parbox{0.91\linewidth}{#1}}\end{center}}

\title{\textbf{The $H_3$ Reflection Group}\\[3pt]
\large Icosahedral construction, spectral formula, and verification}
\author{Computational verification note}
\date{17 July 2026}

\begin{document}
\maketitle

\summarybox{\textbf{Outcome.}  Explicit rotation matrices verify the proposed
$A_4$, $S_4$, and $A_5$ angle formulas in 45 multivariate coefficients each.
The 120 matrices of $W(H_3)$ were reconstructed from the icosahedron; all 139
coefficients through total degree ten agree with the signed class-spectrum
formula, and the one-variable series agrees with
$((1-t^2)(1-t^6)(1-t^{10}))^{-1}$.}

\section{What is being tested}

For a finite matrix group $G\subset GL(V)$, define
\begin{equation}
\Sfun_{G,V}(\mathbf t)=\frac1{|G|}\sum_{g\in G}
\prod_{a\geq1}\det(I-t_a g)^{-1}.                             \tag{1}
\end{equation}
If $g$ has eigenvalues $\lambda_1,\ldots,\lambda_n$, then
\[
\det(I-tg)^{-1}=\sum_{d\geq0}
\chi_{\Sym^dV}(g)t^d.
\]
Thus the coefficient at $t_1^{\tau_1}\cdots t_\ell^{\tau_\ell}$ is
\begin{equation}
\frac1{|G|}\sum_{g\in G}\prod_a h_{\tau_a}(\lambda(g))
=\dim\left(\bigotimes_a\Sym^{\tau_a}V\right)^G.              \tag{2}
\end{equation}
The last equality is the trace of the Reynolds projector.  The test below
constructs the matrices in (2) without using the proposed class data, then
compares against a separate weighted spectral computation.

\section{Floating-point reconstruction diagnostic for $W(H_3)$}

Let $\varphi=(1+\sqrt5)/2$.  Use the twelve vertices
\begin{equation}
\Omega=\{(0,\pm1,\pm\varphi),(\pm1,\pm\varphi,0),
(\pm\varphi,0,\pm1)\}.                                      \tag{3}
\end{equation}
Choose an ordered, linearly independent base triple $B$ of vertices.  For
every ordered target triple $T$ having the same Gram matrix, form the unique
linear map $R=TB^{-1}$.  Retain $R$ when
\[
R^{\mathsf T}R=I,\qquad\det R=1,\qquad R\Omega=\Omega.
\]
This produces 60 distinct rotations.  Adjoining their negatives produces 120
orthogonal transformations.  This is the full symmetry group of the
icosahedron, which is the real reflection group $W(H_3)$.

This ``vertex-frame reconstruction'' is useful when no computer algebra class
table is available, but the present implementation uses floating-point
coordinates and therefore supplies a diagnostic rather than an exact group
certificate.  The implementation reports its equality tolerance, the minimum
separation between reconstructed matrices, the maximum closure-matching
distance, and the condition number of the chosen base triple; it then tests
every one of the $120^2$ products against the reconstructed set.  No closure
claim here is described as exact.

The regenerated diagnostic reports equality tolerance $5\times10^{-10}$,
minimum separation $1.6625$ between distinct reconstructed matrices, maximum
closure-matching distance $3.77\times10^{-16}$, base-triple condition number
$1.852$, and maximum orthogonality residual $4.16\times10^{-16}$.

\section{Polyhedral rotation groups}

For a three-dimensional rotation through angle $\theta$, define
\[
\Phi_\theta(\mathbf t)=\prod_{a\geq1}
\big((1-t_a)(1-2\cos\theta\,t_a+t_a^2)\big)^{-1}.
\]
The spectrum is $(1,e^{i\theta},e^{-i\theta})$, so direct substitution of the
rotation counts gives
\begin{align*}
\Sfun_{A_4,\mathrm{rot}}&=\frac1{12}
(\Phi_0+3\Phi_\pi+8\Phi_{2\pi/3}),\\
\Sfun_{S_4,\mathrm{rot}}&=\frac1{24}
(\Phi_0+6\Phi_{\pi/2}+8\Phi_{2\pi/3}+9\Phi_\pi),\\
\Sfun_{A_5,\mathrm{rot}}&=\frac1{60}
(\Phi_0+15\Phi_\pi+20\Phi_{2\pi/3}
+12\Phi_{2\pi/5}+12\Phi_{4\pi/5}).
\end{align*}
The direct route independently constructs the 12 tetrahedral matrices as the
determinant-one signed permutation matrices preserving a fixed tetrahedron,
the 24 octahedral matrices as all determinant-one signed permutation matrices,
and the 60 icosahedral matrices by the vertex-frame construction above.  For
each group, all 45 coefficients indexed by partitions through total degree
seven agree with the displayed angle formula.

\begin{center}
\begin{tabular}{@{}lrrc@{}}
\toprule
group & order & coefficient checks & result\\
\midrule
$A_4$ tetrahedral rotations & 12 & 45 & \PASS\\
$S_4$ octahedral rotations & 24 & 45 & \PASS\\
$A_5$ icosahedral rotations & 60 & 45 & \PASS\\
\bottomrule
\end{tabular}
\end{center}

\section{The explicit $H_3$ spectral formula}

Write
\[
A_z(\mathbf t)=\prod_{a\geq1}(1-zt_a)^{-1},\qquad
\omega=e^{2\pi i/3},\qquad \xi=e^{2\pi i/5}.
\]
The orientation-preserving subgroup has the following spectra and class sizes:
\begin{center}
\begin{tabular}{@{}clc@{}}
\toprule
size & eigenvalues & geometric type\\
\midrule
1  & $(1,1,1)$ & identity\\
15 & $(1,-1,-1)$ & half turns\\
20 & $(1,\omega,\omega^2)$ & third turns\\
12 & $(1,\xi,\xi^{-1})$ & first fifth-turn class\\
12 & $(1,\xi^2,\xi^{-2})$ & second fifth-turn class\\
\bottomrule
\end{tabular}
\end{center}
Multiplication by central inversion negates all three eigenvalues and supplies
the other five classes.  Grouping (1) by these ten spectra gives
\begin{align}
\Sfun_{H_3}=\frac1{120}\Big[&A_1^3+A_{-1}^3
+15(A_1A_{-1}^2+A_{-1}A_1^2)\notag\\
&+20(A_1A_\omega A_{\omega^2}
+A_{-1}A_{-\omega}A_{-\omega^2})\notag\\
&+12\sum_{j=1}^{2}(A_1A_{\xi^j}A_{\xi^{-j}}
+A_{-1}A_{-\xi^j}A_{-\xi^{-j}})\Big].                       \tag{4}
\end{align}
Equation (4) is exact in all multidegrees.  It is more informative than the
ordinary invariant degrees because mixed coefficients depend on the entire
class spectrum.

\subsection{Ordinary Molien specialization}

Setting only $t_1=t$ nonzero gives
\begin{equation}
\Sfun_{H_3}(t,0,\ldots)=
\frac1{(1-t^2)(1-t^6)(1-t^{10})}.                             \tag{5}
\end{equation}
The right side has degree-$d$ coefficient equal to the number of nonnegative
solutions of $2a+6b+10c=d$.  This elementary count is used as a third route,
independent of both explicit matrix spectra and the multivariate class
calculation.

\section{Verification protocol}

\begin{enumerate}[leftmargin=*,itemsep=3pt]
\item Reconstruct the 60 rotations from (3), adjoin negatives, and test order,
orthogonality, signed determinants, and closure.
\item For every partition $\tau$ through total degree ten, compute the direct
matrix average in (2) over all 120 matrices.
\item Compute the same coefficient using only the ten weighted spectra in (4).
The explicit matrices are not passed to this route.
\item For every $0\leq d\leq10$, compare the direct one-variable average with
the solution count for $2a+6b+10c=d$.
\item Construct the six-dimensional $\Sym^2(V)$ matrices and average them to a
Reynolds projector.
\end{enumerate}

\section{Results}

\begin{center}
\begin{tabularx}{\linewidth}{@{}XrX@{}}
\toprule
Check & observed & status\\
\midrule
number of reconstructed matrices & 120 & \PASS\\
closure products tested & 14,400 & \PASS\\
maximum $\|g^{\mathsf T}g-I\|_F$ & $4.16\times10^{-16}$ & \PASS\\
signed determinant condition & all matrices & \PASS\\
partitions through degree 10 & 139 coefficients & \PASS\\
ordinary Molien degrees $0$ through $10$ & 11 coefficients & \PASS\\
\bottomrule
\end{tabularx}
\end{center}

Some matched multivariate coefficients are
\begin{center}
\begin{tabular}{@{}lrrrrrrr@{}}
\toprule
$\tau$ & $(1)$ & $(2)$ & $(3)$ & $(4)$ & $(2,1)$ & $(2,2)$ & $(3,2)$\\
\midrule
dimension & 0 & 1 & 0 & 1 & 0 & 2 & 0\\
\bottomrule
\end{tabular}
\end{center}
The coefficients of (5) for $d=0,\ldots,10$ are
\[
1,0,1,0,1,0,2,0,2,0,3,
\]
and agree termwise with direct matrix averaging.

For $\tau=(2)$, the target $\Sym^2(V)$ has dimension six.  The explicit
Reynolds projector has
\[
\operatorname{rank}P=1,\qquad \operatorname{tr}P=0.9999999999999998,
\qquad \|P^2-P\|_F=1.26\times10^{-16}.
\]
The invariant is the expected quadratic form.

\section{Reproducibility and limitations}

Run from the repository root:
{\small
\begin{verbatim}
.venv/bin/python -m lambda_distributions.proofs.matrix_group_formula_verification.h3_reflection.verification
\end{verbatim}
}
The companion marimo notebook exposes the maximum tested degree and displays
all coefficient rows.

The group reconstruction uses double-precision values of $\varphi$, not an
exact quadratic-number field.  This does not affect the theoretical proof of
(4), but it makes the group diagnostics numerical.  To control this, vertices
and matrices are canonicalized to nine decimal places, every product is tested,
and all averaged characters must lie within $5\times10^{-7}$ of an integer.
The actual orthogonality and projector residuals are around $10^{-16}$.

The finite checks are diagnostic evidence; the all-degree validity of (4)
comes from the master identity and the complete class spectra.  The computation
specifically verifies that the class weights, signs, fifth-root classes, and
matrix convention have been transcribed consistently.

\end{document}
